\section{Future Work}
\label{sec:ch9-future-work}

\paragraph{Validation against production data.} On the datasets of this work, the probe is more accurate than the coordinator in seven of the eight models, and its disagreement with the coordinator marks coordinator errors, as Section~\ref{subsec:ch8-rq4-answer} states. Every sample of these datasets is generated from a fixed taxonomy of five sub-agents and follows the patterns defined by the generation plan. Requests collected from a coordinator-based multi-agent system in use may not follow a predefined set of patterns, and they may be incomplete, ambiguous, or unrelated to any sub-agent. A labeled dataset of such requests would measure the accuracy of the probe under this variability, and this accuracy would determine how useful the probe is as a misrouting flag in a multi-agent system in use.

A multi-agent system in use may already have a set of production requests annotated with the correct sub-agent, which is used to evaluate the routing decision of its coordinator. Future work could validate the probe on such a system with the following procedure:

\begin{enumerate}
    \item \textbf{Training.} The hidden states of the coordinator are collected by running each annotated request through the prefill step (Section~\ref{subsec:prefill-decode}) of the model in use, with no response token generated. A probe is trained on these hidden states with $k$-fold cross-validation, as for the probes of Section~\ref{subsec:ch7-probe}.
    \item \textbf{Deployment.} The probe is deployed next to the coordinator, on an inference server that returns the hidden states (Section~\ref{subsec:ch9-flag-access}), and it raises the disagreement flag on the incoming requests.
    \item \textbf{Collection period.} The flagged requests are collected for a defined period of time.
    \item \textbf{Annotation.} At the end of the collection period, human annotators examine every flagged request and determine whether the coordinator misrouted it or the flag is a false positive.
\end{enumerate}

The fraction of the flagged requests that are misrouted gives the precision of the flag in real usage. This precision would show how effective a probe trained on evaluation data is as a misrouting flag on the requests of a multi-agent system in use.

\paragraph{Adaptation to a changing sub-agent set.} The sub-agent set of a multi-agent system in use may change over time, for example when a new sub-agent is added or when the responsibilities of an existing one are redefined. Section~\ref{subsec:ch9-flag-cold-start} notes that a probe needs retraining after each of these changes. The labeled data collected before the change may no longer match the new routing task, so the probe could not be used as a misrouting flag until enough new labeled data has accumulated. A change of the sub-agent set also modifies the prompt of the coordinator, because every sub-agent is declared in the prompt as a tool (Section~\ref{subsec:ch7-coordinator-prompt}). The hidden state at the last prompt token is computed from the entire prompt, so it may change even for requests whose correct sub-agent stays the same. The hidden states may also change when the prompt of the coordinator is modified and the sub-agent set stays the same.

Future work could study how the accuracy of an existing probe varies after each type of change: the addition, the removal, and the redefinition of a sub-agent. This study would show which changes leave the accuracy of the probe unaffected and which ones require new labeled data. Future work could also study whether an existing probe can be adapted to the new sub-agent set by reusing what it has already learned, for example by updating it with only the samples labeled after the change. If the adapted probe needs fewer labeled samples than a new probe to reach the same accuracy, fewer requests would need to be labeled after each change before the probe can be used again as a misrouting flag.

\paragraph{Detailed analysis of multi-turn behavior.} On the multi-turn dataset, the multi-turn drop of the probe is larger than the multi-turn drop of the coordinator in all seven models, as Section~\ref{subsec:ch8-multi-turn} reports. Section~\ref{subsec:ch9-limitations-multiturn} states that every conversation in this dataset has the same structure and the same number of messages. A conversation in a multi-agent system in use may span many more turns, change topic, and include turns unrelated to the routing decision. The hidden state at the last prompt token is computed from the entire conversation, so the information needed for routing may be spread across more turns or placed far from the last prompt token. The accuracy of the probe on such conversations may therefore differ from the one measured on the multi-turn dataset of this work.

Future work could build a synthetic multi-turn dataset in which the number of turns and the structure of the conversation vary as design factors. The structure could follow a defined taxonomy of patterns, such as topic drift, turns unrelated to the routing decision, and different positions of the turn that carries the routing information. Such a dataset would allow a study of how the accuracy of the probe varies across the patterns, and across shorter and longer conversations. The same dataset could also be used to test whether a probe that reads the hidden states at several positions, such as the last token of each turn, is more accurate on long conversations, or on conversations with topic drift, than a probe that reads only the last prompt token. The results would show whether the probe remains usable as a misrouting flag on conversations longer and less regular than the ones of this work.
